\section{Ordinary quadratic indicators and fixed-shift pairs}
\label{sec:quadratic-pairs}

The same comparison also treats ordinary quadratic norm indicators.
In this application inert primes may occur to an even exponent; they
are not excluded altogether. The rough Dirichlet series still has
two principal character branches, and their projection retains the
conditions at ramified primes. This gives a common statement for
the five pair conjectures considered by Smilansky
\cite[Sections 6.1 and 6.3]{smilansky}.

\subsection{A quadratic character and finitely many valuation rules}

Let \(\chi\) be a primitive nonprincipal quadratic Dirichlet
character of conductor \(q\). Fix a finite set of primes
\(B_0\) containing every prime divisor of \(q\), and masks
\(\alpha_p:\Z_{\ge0}\to\{0,1\}\), \(p\in B_0\), with
\(\alpha_p(0)=1\). Define the multiplicative indicator
\[
 f(n)=\prod_{p\in B_0}\alpha_p(v_p(n))
       \prod_{\substack{p\notin B_0\\\chi(p)=-1}}
                                      \1_{2\mid v_p(n)}
 \quad(n\ge1),\qquad f(n)=0\quad(n\le0).
\]
Thus split primes outside \(B_0\) have no valuation restriction.
The finite masks are fixed; no periodicity in their exponent is
required. Write \(F_p(s)=\sum_{e\ge0}\alpha_p(e)p^{-es}\)
for \(p\in B_0\).

The bad-prime factor must retain a joint residue condition. Let
\(x=(x_p)_{p\in B_0}\in\prod_{p\in B_0}\Z_p\), and put
\(e_p=v_p(x_p)\) when all these valuations are finite. For each
\(\ell\mid q\), define a unit modulo \(\ell^{v_\ell(q)}\) by
\[
 u_\ell(x)=\ell^{-e_\ell}x_\ell
            \prod_{\substack{p\in B_0\\p\ne\ell}}p^{-e_p}
                                  \pmod{\ell^{v_\ell(q)}}.
\]
These residues define \(u_{B_0}(x)\in(\Z/q\Z)^\times\) by CRT.
At an ordinary positive integer this is exactly
\(n/n_{B_0}\pmod q\). Set
\begin{align}
 \vartheta_{B_0}(x)
   &=\prod_{p\in B_0}\alpha_p(e_p)
                    \1_{\chi(u_{B_0}(x))=1},\notag\\
 W_{B_0}(h)
   &=\int_{\prod_{p\in B_0}\Z_p}
             \vartheta_{B_0}(x)\vartheta_{B_0}(x+h)\,dx.
                                                        \label{eq:quadratic-bad-mass}
\end{align}
Normalized product Haar measure is used. Define the integrand to
be zero on the null sets with an infinite valuation. In general
\(W_{B_0}\) is coupled: it is not a product of separate
primewise tests.

\begin{theorem}\label{thm:quadratic-indicator-pairs}
Assume the analytic inputs of Section~\ref{sec:inputs}. For the
fixed data above and every fixed \(h\ne0\),
\begin{equation}\label{eq:quadratic-pair-asymptotic}
 \sum_{1\le n\le X}f(n)f(n+h)
       =C_{\chi,\boldsymbol\alpha}(h)\frac X{\log X}
                            +o_{\chi,\boldsymbol\alpha,h}
                                      \left(\frac X{\log X}\right),
\end{equation}
where
\begin{align}
 C_{\chi,\boldsymbol\alpha}(h)
 &=\frac{4L(1,\chi)}\pi\,W_{B_0}(h)
       \prod_{p\in B_0}\frac{1-\chi(p)/p}{1-1/p}\notag\\
 &\quad\cdot
       \prod_{\substack{p\notin B_0,\ \chi(p)=-1\\p\mid h}}
                     \frac{1-p^{-v_p(h)-1}}{1-p^{-1}}.
                                                        \label{eq:quadratic-pair-constant}
\end{align}
If \(W_{B_0}(h)=0\), the pair count is identically zero.
Otherwise the constant is positive. The theorem concerns the
specified Euler indicator; it does not identify that indicator
with the values of an arbitrary principal quadratic form.
\end{theorem}

\begin{proof}
We verify the four hypotheses of
Theorem~\ref{thm:model-comparison}, then evaluate the model product.
Choose a fixed even \(L_0\) divisible by \(q\) and all primes
in \(B_0\), and use the prime class one modulo \(L_0\).
Put \(L=\log X\), \(y=e^{L^{2/5}}\), \(S=\{p:p\le y\}\),
\(E=e^{\sqrt L}\), and \(V_S=\prod_{p\le y}(1-1/p)\).

The exact Euler identity for \(F(s)=\sum f(n)n^{-s}\) is
\begin{equation}\label{eq:quadratic-euler-square}
 F(s)^2=\zeta(s)L(s,\chi)
 \prod_{\substack{p\notin B_0\\\chi(p)=-1}}(1-p^{-2s})^{-1}
 \prod_{p\in B_0}F_p(s)^2(1-p^{-s})(1-\chi(p)p^{-s}).
\end{equation}
The first correction converges absolutely for \(\Re s>1/2\).
Every finite factor is analytic near one and positive there on
the real axis. With this branch convention, the one-variable
constant \(c_f\) satisfies
\begin{equation}\label{eq:quadratic-marginal}
 c_f^2=\frac{L(1,\chi)}\pi
 \prod_{\substack{p\notin B_0\\\chi(p)=-1}}(1-p^{-2})^{-1}
 \prod_{p\in B_0}F_p(1)^2(1-p^{-1})(1-\chi(p)p^{-1}).
\end{equation}
Here \(F_p(1)\) denotes the local factor evaluated at \(s=1\),
not at its series variable. The factor \(\pi=\Gamma(1/2)^2\)
is the Selberg--Delange normalization.

Remove all factors at primes in \(S\), and write the resulting
rough series as \(F_R(s)=(s-1)^{-1/2}H_R(s)\). Its leading
coefficient is
\[
 H_R(1)/\Gamma(1/2)
          =c_f\prod_{p\le y}F_{f,p}(1)^{-1},
\]
where \(F_{f,p}\) is the actual local factor of \(f\).
The Cauchy bound
\(H_R^{(j)}(1)/j!\ll_j H_R(1)(C\log y)^j\) permits the
fixed-depth density
\[
 M_K(t)=\sum_{j=0}^K\frac{H_R^{(j)}(1)}{j!}
              \frac{(\log t)^{-1/2-j}}{\Gamma(1/2-j)}.
\]
Use this formula throughout \(X^{1/4}\le t\le X^3\), and
extend it smoothly and nonnegatively below \(X^{1/4}\) with its
leading upper bound. This contains the rough range in (C3), since
\(Y/r\ge X^{1/2}/E>X^{1/4}\) for large \(X\).
For \(n=av\), \(a=n_S\), define
\begin{equation}\label{eq:quadratic-model}
 D(n)=f(a)(1+\chi(v))M_K(v)/V_S.
\end{equation}
For every fixed depth, \(D\ll L^{-3/10}\) on the relevant
polynomial ranges, so (C1) holds eventually with a depth-independent
bound one. Multiplication by a split bank prime changes only
\(a\), preserving \(f(a)\) and \(v\) exactly. This proves (C2),
including repeated prime factors.

For (C3), inflate a test character \(\eta\) to a modulus
containing \(q\). The square of its rough Dirichlet series has
factors \(L(s,\eta)L(s,\eta\chi)\), an absolutely convergent
prime-square correction, and the removed small-prime factors.
All primes dividing the polylogarithmic test modulus lie in \(S\).
Only \(\eta=1\) and \(\eta=\chi\) give a principal branch,
and their rough coefficients agree pointwise. The factor
\(1+\chi(v)\) in \eqref{eq:quadratic-model} reproduces both.
The proof of Lemma~\ref{lem:rough-means} now applies with exponent
\(1/2\): the contour is \(1-c/\log y\), its height is
\(e^{L^{1/10}}\), and the Taylor remainder has relative size
\(O_K(L^{-3(K+1)/5})\). The inverted inert-prime correction
in \eqref{eq:quadratic-euler-square} has the same absolute
convergence and derivative bounds as the correction in that lemma.
The finite masked factors have been removed completely from the
rough variable. Thus they introduce no further rough branches.
Conditioning on \(a\le E\) leaves only its fixed multiplier
\(f(a)\), outside the periodic rough sum. Choosing \(K\)
after the desired logarithmic accuracy proves (C3).

For (C4), keep exact bank records \(r,s\le E\); larger ones
have contribution \(XL^{O(1)}e^{-cL^{1/10}}\) by
\eqref{eq:record-tail}. The retained records are coprime because
bank primes dividing \(h\) have been deleted. Their divisibilities
give one progression with \(|J|/(rs)+O(1)\) terms.
Exactness contributes two excluded bank roots outside \(rs\),
and one at a prime of \(rs\). Together with \(1/(rs)\), these
are exactly the record masses \eqref{eq:bank-law}.

At an inert prime outside \(B_0\), a necessary condition is
\(v_p(n)\ne1\). It excludes \(p-1\) classes modulo \(p^2\).
For \(p\nmid h\), the two endpoints exclude \(2(p-1)\)
classes; bank-record slopes are units at these primes. Drop the
finitely many primes dividing \(2h\prod_{p\in B_0}p\) for
an upper bound. The nonbank sieve dimension is one, with a
convergent \(1+O(p^{-2})\) correction to the ordinary two-root
product. A Selberg upper sieve with prime-square atoms, level
\(z=X^{1/100}\), and \(w=z^{1/16}\), has endpoint error
\(O(z^{4+o(1)})\) per retained record. This follows by summing
the CRT root errors: a least-common-multiple atom has at most
\(2^{\omega(d)}d\) roots and modulus \(d^2\).
Summing the \(E^2=X^{o(1)}\) records remains power smaller
than \(X^{1/2}\).

For \(ff\), \(Df\), and \(DD\), the nonbank sieve
denominators have orders respectively
\(\log w\), \((\log y\log w)^{1/2}\), and \(\log y\).
The bound \(D\ll L^{-3/10}\) supplies the missing factor
on each model side. Explicitly, since \(\log y=L^{2/5}\)
and \(\log w\asymp L\), their factors are
\[
 (\log w)^{-1}\ll L^{-1},\qquad
 \frac{L^{-3/10}}{(\log y\log w)^{1/2}}\ll L^{-1},\qquad
 \frac{L^{-3/5}}{\log y}=L^{-1}.
\]
All three cases therefore give the common
scale \(\mathcal S_X=X/L\), with the exact bank masses.
This proves (C4), uniformly in the interval \(J\). The exponents
in the exceptional-record errors do not depend on the final
fixed depth. The comparison theorem now gives
\begin{equation}\label{eq:quadratic-model-comparison}
 \sum_{X\le n\le2X}\{f(n)f(n+h)-D(n)D(n+h)\}
                                      =o_h(X/L).
\end{equation}

It remains to calculate the model main. Let
\(I_y=\{p\le y:p\notin B_0,\ \chi(p)=-1\}\) and put
\[
 \Lambda_y=\frac{2c_f}{\sqrt L}
      \prod_{p\in B_0}\frac{F_p(1)^{-1}}{1-p^{-1}}
      \prod_{p\in I_y}(1+p^{-1}).
\]
Outside the negligible records \(n_S>E\), the rough cofactor
has \(\log v=L+O(\sqrt L)\). If the small inert exponents
are even, its character is \(\chi(n/n_{B_0})\). Hence
\begin{equation}\label{eq:quadratic-model-leading}
 D(n)=\Lambda_y\vartheta_{B_0}(n)
            \prod_{p\in I_y}\1_{2\mid v_p(n)}
                                      (1+O_K(L^{-1/2})).
\end{equation}
At an inert prime the local pair density is
\begin{equation}\label{eq:quadratic-inert-pair}
 \delta_p(h)=\frac p{p+1}(1-p^{-v_p(h)-1}).
\end{equation}
To see this, put \(e=v_p(h)\) and partition \(v_p(n)\)
below, equal to, and above \(e\). Below \(e\), both endpoints
have the same valuation, and the permitted even terms form a
geometric sum. If \(e\) is odd there are no further terms.
If \(e\) is even, the stratum with both valuations equal to
\(e\) has mass \(p^{-e}(p-2)/p\); the two possible higher
even tails have combined mass \(2p^{-e}/(p(p+1))\).
The two parity cases give \eqref{eq:quadratic-inert-pair}, also
for \(p=2\).

We justify the product of these densities at the growing cutoff.
Extract the inert square parts \(d^2,e^2\), with \(d,e\)
supported on \(I_y\), and retain \(d,e\le E\). Truncate
each bad-prime power at \(E\). The omitted integer count is
\(O_{B_0,h}(X/E+\sqrt X)\), and the omitted Haar mass is
\(O_{B_0,h}(E^{-1})\). Each retained joint record and residual
unit state gives an empty stratum or finitely many progressions
modulo \(X^{o(1)}\). On them the two residual forms must be
units at the primes in \(I_y\). Empty root cells are discarded;
the remaining sieve dimension is at most one. There are at most
\(O_{B_0,q}(E^2(\log E)^{2|B_0|})\) retained cells: two
square prefixes, two bounded valuation lists, and a fixed number
of residual unit choices. Their moduli are at most
\(O_{B_0,q}(E^{4+2|B_0|})\). Use upper and lower
fundamental-lemma weights of level \(D_0=X^{1/100}\).
Their parameter \(\log D_0/\log y\asymp L^{3/5}\)
gives relative error smaller than every log power. Since the
root count at a squarefree sieve divisor is at most
\(2^{\omega(d)}\), a conservative total endpoint bound is
\[
 E^{O_{B_0,q}(1)}D_0^2(\log D_0)^{O(1)}
                         =X^{1/50+o(1)}.
\]
This is negligible compared with \(X/L^A\) for every fixed
\(A\). Their main masses are summable: compatibility requires
\((d,e)^2\mid h\), and their progression densities are
\(\ll_h1/\operatorname{lcm}(d^2,e^2)\), with geometric
bad-prime tails. Summing the disjoint valuation cells proves
\[
 \sum_{X\le n\le2X}D(n)D(n+h)
   =\Lambda_y^2XW_{B_0}(h)\prod_{p\in I_y}\delta_p(h)
                                                 +o_h(X/L).
\]
For \(p\nmid h\),
\((1+p^{-1})^2\delta_p(h)=1-p^{-2}\); for \(p\mid h\)
the additional ratio is
\((1-p^{-v_p(h)-1})/(1-p^{-1})\). Substitution of
\eqref{eq:quadratic-marginal} cancels every \(F_p(1)\)
and the inert Euler product, leaving
\eqref{eq:quadratic-pair-constant}. This proves the dyadic
asymptotic. Summing dyadically gives
\eqref{eq:quadratic-pair-asymptotic}; the upper bound from (C4)
controls the part below \(\varepsilon X\), after which
\(\varepsilon\) tends to zero.

Finally, if an actual positive pair has \(f(n)f(n+h)=1\),
then both residual characters equal one. Its valuations and
residual units are preserved on a sufficiently deep joint
\(B_0\)-adic cylinder. That cylinder has positive Haar measure,
so \(W_{B_0}(h)>0\). This proves the exact zero assertion.
\end{proof}

\subsection{The five quadratic forms}

Let \(b_d(n)\) indicate the positive integers of the form
\(x^2+dy^2\), with \(x,y\in\Z\). Coordinates may vanish,
and no coprimality condition is imposed. Put \(b_d(n)=0\)
for \(n\le0\). Smilansky's
prime-factor criterion \cite[Theorem 10]{smilansky} identifies
these ordinary indicators with the preceding Euler indicators:
\[
\begin{array}{c|c|c|c}
d&\chi&B_0&\text{nontrivial mask in }B_0\\\hline
1&\chi_{-4}&\{2\}&\text{none}\\
2&\chi_{-8}&\{2\}&\text{none}\\
3&\chi_{-3}&\{3\}&\text{none}\\
4&\chi_{-4}&\{2\}&\alpha_2(e)=\1_{e\ne1}\\
7&\chi_{-7}&\{2,7\}&\alpha_2(e)=\1_{e\ne1}
\end{array}
\]
The corresponding order discriminants are \(-4,-8,-12,-16,-28\).
In particular the third, fourth and fifth forms are conductor-two
orders. For \(d=3\), the ordinary value set nevertheless has
exactly the inert-even rule for \(\chi_{-3}\), including the
prime two; no extra dyadic mask is imposed.

\begin{corollary}\label{cor:smilansky-five-pairs}
Under the inputs of Section~\ref{sec:inputs}, for
\(d\in\{1,2,3,4,7\}\) and fixed \(h\ne0\),
\[
 \sum_{1\le n\le X}b_d(n)b_d(n+h)
       =\mathscr T_{d,h}\frac X{\log X}+o_{d,h}(X/\log X),
\]
with precisely the coefficients predicted in
\cite[Section 6.3, equations (6.2)--(6.4)]{smilansky}.
In particular,
\begin{equation}\label{eq:minus-eight-shift-one}
 \sum_{n\le X}b_2(n)b_2(n+1)
                       \sim\frac1{\sqrt2}\frac X{\log X}.
\end{equation}
\end{corollary}

\begin{proof}
Write \(Q_d=\{p:\chi(p)=-1\}\), except that no prime in
\(B_0\) is included; for the displayed five forms this exception
changes none of the inert sets. These are respectively
\[
 \{p\equiv3\ (4)\},\quad \{p\equiv5,7\ (8)\},\quad
 \{p\equiv2\ (3)\},\quad \{p\equiv3\ (4)\},\quad
 \{p\equiv3,5,6\ (7)\}.
\]
Let \(R_1=R_2=R_4=\{2\}\), \(R_3=\{3\}\), and
\(R_7=\{2,7\}\). The bad-prime integrand factors in these
cases. For \(d=7\), this uses \(\chi_{-7}(2)=1\), which
removes the cross-valuation phase from its ramified unit test.
Put \(a_p=v_p(|h|)\). The local densities are
\begin{align*}
 W_{1,2}(h)&=\begin{cases}
  1/4,&a_2=0,\\ 1/2-3/2^{a_2+2},&a_2\ge1,
 \end{cases}
 &W_{2,2}(h)&=\begin{cases}
  1/4,&a_2=0,1,\\ 1/2-3/2^{a_2+1},&a_2\ge2,
 \end{cases}\\[2mm]
 W_{3,3}(h)&=\tfrac12(1-2/3^{a_3+1}),
 &W_{7,7}(h)&=\tfrac12(1-4/7^{a_7+1}),\\[2mm]
 W_{4,2}(h)&=\begin{cases}
  1/8,&a_2=0,\\0,&a_2=1,\\5/16,&a_2=2,\\
  3/8-3/2^{a_2+2},&a_2\ge3,
 \end{cases}
 &W_{7,2}(h)&=\begin{cases}
  1/2,&a_2=0,1,\\3/4,&a_2\ge2.
 \end{cases}
\end{align*}
For completeness, the odd ramified unit predicate is a square
modulo \(p=3\) or \(7\). At \(a_p=0\), there are
\((p-3)/4\) nonzero square pairs modulo \(p\), and exactly
one zero neighborhood contributes \(1/(2p)\). The result is
\((p-1)/(4p)\). For \(a_p\ge1\), the recurrence
\(W(a_p)=(p-1)/(2p)+W(a_p-1)/p\) gives the displayed formulas.
At two the ramified unit predicates are respectively one modulo
four and one or three modulo eight. Their initial pair masses
are \(1/4\), and, for the latter, \(1/4\) also when \(a_2=1\).
The former has mass \(1/8\) at \(a_2=1\); thereafter its
recurrence is \(W(a)=1/4+W(a-1)/2\). The latter has mass
\(1/8\) at \(a_2=2\) and obeys that recurrence for \(a\ge3\).
Removing valuation one gives the \(d=4\) row: for \(a\ge2\)
its mass is \(1/4+W_{1,2}(h/4)/4\). The \(d=7\) dyadic
row counts only the two valuation masks, giving its stated masses.

Using primitive characters, the prefactors are
\[
 c_d=\frac{4L(1,\chi)}\pi
                 \prod_{p\mid\operatorname{cond}(\chi)}(1-p^{-1})^{-1},
 \qquad
 (c_1,c_2,c_3,c_4,c_7)
       =(2,2\sqrt2,2/\sqrt3,2,2\sqrt7/3).
\]
For \(d=7\) the extra bad prime two contributes ratio one
in \eqref{eq:quadratic-pair-constant}. Consequently
\[
 \mathscr T_{d,h}=c_d\prod_{p\in R_d}W_{d,p}(h)
        \prod_{\substack{p\in Q_d\\p\mid h}}
                    \frac{1-p^{-a_p-1}}{1-p^{-1}},
\]
as required. The primitive values are
\(L(1,\chi_{-3})=\pi/(3\sqrt3)\) and
\(L(1,\chi_{-7})=\pi/\sqrt7\); Smilansky's series in
his equation (6.3) omits the Euler factor at two.
At \(h=1\) the five constants are
\[
 (\mathscr T_{1,1},\mathscr T_{2,1},\mathscr T_{3,1},
      \mathscr T_{4,1},\mathscr T_{7,1})
   =\left(\frac12,\frac1{\sqrt2},\frac1{3\sqrt3},
                        \frac14,\frac1{2\sqrt7}\right).
\]
Only \(d=4\) has a zero coefficient for nonzero shifts:
\(v_2(h)=1\). Its pair count is exactly zero because its
values are zero or one modulo four. All other displayed
coefficients are positive.
\end{proof}
